% TOP_PROBLEM: {"problemId": "problem.list-edge-coloring-conjecture", "canonicalTitle": "List edge-coloring conjecture", "releaseRank": 496, "primaryDomain": "combinatorics_discrete_geometry", "primaryDomainLabel": "Combinatorics & discrete geometry", "displayStatus": "open_with_solved_subcases", "releaseStatus": "open_with_solved_subcases"}
% !TeX program = lualatex
% Definition and sources only; this file is not an LLM solution attempt.
\documentclass[11pt]{article}
\usepackage[margin=25mm]{geometry}
\usepackage{fontspec}
\setmainfont{Latin Modern Roman}
\usepackage{amsmath,amssymb,mathtools}
\usepackage{unicode-math}
\setmathfont{Latin Modern Math}
\usepackage{xurl,hyperref}
\hypersetup{colorlinks=true,urlcolor=blue}
\setcounter{secnumdepth}{0}
\setlength{\parindent}{0pt}
\setlength{\parskip}{5pt}
\setlength{\emergencystretch}{3em}
\begin{document}
\section*{\texorpdfstring{List edge-coloring conjecture}{List edge-coloring conjecture}}\label{496}
\subsection{Definitions and mathematical statement}
Let $G$ be a finite loopless multigraph, allowing parallel edges. A proper edge coloring assigns different colors to every pair of edges sharing a vertex. The ordinary chromatic index $\chi'(G)$ is the smallest number of colors permitting one. Its list chromatic index $\chi'_\ell(G)$ is the least $k$ such that, for every assignment of a list $L(e)$ of $k$ permitted colors to each edge, a proper edge coloring $c$ exists with $c(e)\in L(e)$.

The conjecture is
\[
 \chi'_\ell(G)=\chi'(G)\qquad\text{for every finite loopless multigraph }G.
\]
Different edges may have entirely different lists; taking all lists equal gives only the ordinary coloring problem. The statement concerns edge lists, not vertex-list coloring, and the absence of loops is essential to the standard proper-coloring definition.
\par\smallskip\textit{Formulation and concept references:} \href{https://www.sciencedirect.com/science/article/pii/S0012365X25000287}{[S1]}.

\subsection{Short English statement}
Does assigning each edge its own list of colors ever require larger lists than the number of colors needed for ordinary edge coloring?
\subsection{Sources}
\begin{itemize}
\item[S1]Multigraph edge-coloring with local list sizes.
\url{https://www.sciencedirect.com/science/article/pii/S0012365X25000287}.
\textit{Formulation source.}
\end{itemize}
\end{document}
